\section{A counterexample}

Let $N=\{1,2,3,4,5\}$ and let the universe of features be
$\{A,B,C,D\}$.  Associate to the five elements the feature sets
\[
\begin{aligned}
 E_1&=\{A,B\}, & E_2&=\{A,B,C,D\}, & E_3&=\{C,D\},\\
 E_4&=\{A,C\}, & E_5&=\{B,D\}.&&
\end{aligned}
\]
For $S\subseteq N$, define
\[
  f(S)=\left|\bigcup_{i\in S}E_i\right|.
\]
This coverage function is nonnegative, monotone, and submodular
\cite{bach2013}.  Fix the marginal vector
\[
  (x_1,x_2,x_3,x_4,x_5)
  =\left(\frac3{10},\frac7{20},\frac3{10},\frac7{20},\frac7{20}\right).
\]

\subsection{The unrestricted optimum}

Consider the following three-point distribution.
\begin{table}[h]
\centering
\caption{Primal feasible and optimal distribution for $f^+(\boldsymbol{x})$.}
\label{tab:primal}
\begin{tabular}{cc}
$S$ & $\theta(S)$\\ \hline
$\{1,3\}$ & $3/10$\\
$\{2\}$ & $7/20$\\
$\{4,5\}$ & $7/20$\\
all other sets & $0$
\end{tabular}
\end{table}
Its weights sum to one and have the prescribed marginals.  Each of its three
support sets covers all four features, so its expected value is $4$.  Since
$f(S)\leq4$ for every $S\subseteq N$, this proves
\[
  f^+(\boldsymbol{x})=4.
\]

\subsection{A pairwise-independent dual certificate}

The dual of \eqref{eq:pairwise} is
\[
\begin{aligned}
 f^{++}(\boldsymbol{x})=\min\quad
 &\sum_{i<j}\lambda_{ij}x_i x_j+
   \sum_i\lambda_i x_i+\lambda_0\\
 \text{subject to}\quad
 &\sum_{\substack{i<j\\i,j\in S}}\lambda_{ij}
   +\sum_{i\in S}\lambda_i+\lambda_0\geq f(S)
   \qquad(S\subseteq N).
\end{aligned}
\]
Write the left side of the constraint indexed by $S$ as $\lambda(S)$.  Take
\[
\begin{gathered}
 \lambda_0=\frac12,\qquad
 \lambda_1=\lambda_3=\lambda_4=\lambda_5=\frac32,\qquad
 \lambda_2=\frac72,\\
 \lambda_{12}=\lambda_{23}=-1,\qquad
 \lambda_{24}=\lambda_{25}=-\frac32,\qquad
 \lambda_{13}=\lambda_{45}=\frac12,\\
 \lambda_{14}=\lambda_{15}=\lambda_{34}=\lambda_{35}=-\frac12.
\end{gathered}
\]
The following table verifies every one of the $2^5=32$ dual constraints.
\begin{longtable}{ccc}
\caption{Dual feasible solution for $f^{++}(\boldsymbol{x})$.}
\label{tab:dual}
\\
$S$ & $\lambda(S)$ & $f(S)$\\ \hline
$\emptyset$ & $1/2$ & $0$\\
$\{1\}$ & $2$ & $2$\\
$\{2\}$ & $4$ & $4$\\
$\{3\}$ & $2$ & $2$\\
$\{4\}$ & $2$ & $2$\\
$\{5\}$ & $2$ & $2$\\
$\{1,2\}$ & $9/2$ & $4$\\
$\{1,3\}$ & $4$ & $4$\\
$\{1,4\}$ & $3$ & $3$\\
$\{1,5\}$ & $3$ & $3$\\
$\{2,3\}$ & $9/2$ & $4$\\
$\{2,4\}$ & $4$ & $4$\\
$\{2,5\}$ & $4$ & $4$\\
$\{3,4\}$ & $3$ & $3$\\
$\{3,5\}$ & $3$ & $3$\\
$\{4,5\}$ & $4$ & $4$\\
$\{1,2,3\}$ & $11/2$ & $4$\\
$\{1,2,4\}$ & $4$ & $4$\\
$\{1,2,5\}$ & $4$ & $4$\\
$\{1,3,4\}$ & $9/2$ & $4$\\
$\{1,3,5\}$ & $9/2$ & $4$\\
$\{1,4,5\}$ & $9/2$ & $4$\\
$\{2,3,4\}$ & $4$ & $4$\\
$\{2,3,5\}$ & $4$ & $4$\\
$\{2,4,5\}$ & $9/2$ & $4$\\
$\{3,4,5\}$ & $9/2$ & $4$\\
$\{1,2,3,4\}$ & $9/2$ & $4$\\
$\{1,2,3,5\}$ & $9/2$ & $4$\\
$\{1,2,4,5\}$ & $4$ & $4$\\
$\{1,3,4,5\}$ & $11/2$ & $4$\\
$\{2,3,4,5\}$ & $4$ & $4$\\
$\{1,2,3,4,5\}$ & $4$ & $4$\\
\end{longtable}

For any pairwise-independent feasible distribution, taking expectations of
the pointwise inequality $f(S)\leq\lambda(S)$ and using its prescribed first
and second moments gives
\[
\begin{aligned}
\E[\lambda(S)]={}&-(x_1x_2+x_2x_3)
-\frac32(x_2x_4+x_2x_5)+\frac12(x_1x_3+x_4x_5)\\
&-\frac12(x_1x_4+x_1x_5+x_3x_4+x_3x_5)\\
&+\frac32(x_1+x_3+x_4+x_5)+\frac72x_2+\frac12
=\frac{479}{160}.
\end{aligned}
\]
Thus weak duality yields
\[
  f^{++}(\boldsymbol{x})\leq\frac{479}{160}.
\]
Combining this with the unrestricted optimum gives
\[
 \frac{f^+(\boldsymbol{x})}{f^{++}(\boldsymbol{x})}
 \geq\frac{4}{479/160}=\frac{640}{479}>\frac43.
\]
Therefore the proposed $4/3$ upper bound fails for this instance.
